% Options for packages loaded elsewhere
% Options for packages loaded elsewhere
\PassOptionsToPackage{unicode}{hyperref}
\PassOptionsToPackage{hyphens}{url}
\PassOptionsToPackage{dvipsnames,svgnames,x11names}{xcolor}
%
\documentclass[
  spanish,
  10pt,
  letterpaper,
  twocolumn]{article}
\usepackage{xcolor}
\usepackage[top=1.0cm,bottom=1.5cm,left=1.5cm,right=1.5cm]{geometry}
\usepackage{amsmath,amssymb}
\setcounter{secnumdepth}{-\maxdimen} % remove section numbering
\usepackage{iftex}
\ifPDFTeX
  \usepackage[T1]{fontenc}
  \usepackage[utf8]{inputenc}
  \usepackage{textcomp} % provide euro and other symbols
\else % if luatex or xetex
  \usepackage{unicode-math} % this also loads fontspec
  \defaultfontfeatures{Scale=MatchLowercase}
  \defaultfontfeatures[\rmfamily]{Ligatures=TeX,Scale=1}
\fi
\usepackage{lmodern}
\ifPDFTeX\else
  % xetex/luatex font selection
  \setmainfont[]{Fira Sans}
\fi
% Use upquote if available, for straight quotes in verbatim environments
\IfFileExists{upquote.sty}{\usepackage{upquote}}{}
\IfFileExists{microtype.sty}{% use microtype if available
  \usepackage[]{microtype}
  \UseMicrotypeSet[protrusion]{basicmath} % disable protrusion for tt fonts
}{}
\usepackage{setspace}
\makeatletter
\@ifundefined{KOMAClassName}{% if non-KOMA class
  \IfFileExists{parskip.sty}{%
    \usepackage{parskip}
  }{% else
    \setlength{\parindent}{0pt}
    \setlength{\parskip}{6pt plus 2pt minus 1pt}}
}{% if KOMA class
  \KOMAoptions{parskip=half}}
\makeatother
% Make \paragraph and \subparagraph free-standing
\makeatletter
\ifx\paragraph\undefined\else
  \let\oldparagraph\paragraph
  \renewcommand{\paragraph}{
    \@ifstar
      \xxxParagraphStar
      \xxxParagraphNoStar
  }
  \newcommand{\xxxParagraphStar}[1]{\oldparagraph*{#1}\mbox{}}
  \newcommand{\xxxParagraphNoStar}[1]{\oldparagraph{#1}\mbox{}}
\fi
\ifx\subparagraph\undefined\else
  \let\oldsubparagraph\subparagraph
  \renewcommand{\subparagraph}{
    \@ifstar
      \xxxSubParagraphStar
      \xxxSubParagraphNoStar
  }
  \newcommand{\xxxSubParagraphStar}[1]{\oldsubparagraph*{#1}\mbox{}}
  \newcommand{\xxxSubParagraphNoStar}[1]{\oldsubparagraph{#1}\mbox{}}
\fi
\makeatother


\usepackage{longtable,booktabs,array}
\newcounter{none} % for unnumbered tables
\usepackage{calc} % for calculating minipage widths
% Correct order of tables after \paragraph or \subparagraph
\usepackage{etoolbox}
\makeatletter
\patchcmd\longtable{\par}{\if@noskipsec\mbox{}\fi\par}{}{}
\makeatother
% Allow footnotes in longtable head/foot
\IfFileExists{footnotehyper.sty}{\usepackage{footnotehyper}}{\usepackage{footnote}}
\makesavenoteenv{longtable}
\usepackage{graphicx}
\makeatletter
\newsavebox\pandoc@box
\newcommand*\pandocbounded[1]{% scales image to fit in text height/width
  \sbox\pandoc@box{#1}%
  \Gscale@div\@tempa{\textheight}{\dimexpr\ht\pandoc@box+\dp\pandoc@box\relax}%
  \Gscale@div\@tempb{\linewidth}{\wd\pandoc@box}%
  \ifdim\@tempb\p@<\@tempa\p@\let\@tempa\@tempb\fi% select the smaller of both
  \ifdim\@tempa\p@<\p@\scalebox{\@tempa}{\usebox\pandoc@box}%
  \else\usebox{\pandoc@box}%
  \fi%
}
% Set default figure placement to htbp
\def\fps@figure{htbp}
\makeatother



\ifLuaTeX
\usepackage[bidi=basic,shorthands=off]{babel}
\else
\usepackage[bidi=default,shorthands=off]{babel}
\fi
\ifPDFTeX
\else
\babelfont{rm}[]{Fira Sans}
\fi
\ifLuaTeX
  \usepackage{selnolig} % disable illegal ligatures
\fi


\setlength{\emergencystretch}{3em} % prevent overfull lines

\providecommand{\tightlist}{%
  \setlength{\itemsep}{0pt}\setlength{\parskip}{0pt}}



 


\usepackage{tikz}
\usepackage{graphicx}
\usepackage{amsmath}
\renewcommand{\arraystretch}{1.4}
\makeatletter
\@ifpackageloaded{caption}{}{\usepackage{caption}}
\AtBeginDocument{%
\ifdefined\contentsname
  \renewcommand*\contentsname{Tabla de contenidos}
\else
  \newcommand\contentsname{Tabla de contenidos}
\fi
\ifdefined\listfigurename
  \renewcommand*\listfigurename{Listado de Figuras}
\else
  \newcommand\listfigurename{Listado de Figuras}
\fi
\ifdefined\listtablename
  \renewcommand*\listtablename{Listado de Tablas}
\else
  \newcommand\listtablename{Listado de Tablas}
\fi
\ifdefined\figurename
  \renewcommand*\figurename{Figura}
\else
  \newcommand\figurename{Figura}
\fi
\ifdefined\tablename
  \renewcommand*\tablename{Tabla}
\else
  \newcommand\tablename{Tabla}
\fi
}
\@ifpackageloaded{float}{}{\usepackage{float}}
\floatstyle{ruled}
\@ifundefined{c@chapter}{\newfloat{codelisting}{h}{lop}}{\newfloat{codelisting}{h}{lop}[chapter]}
\floatname{codelisting}{Listado}
\newcommand*\listoflistings{\listof{codelisting}{Listado de Listados}}
\makeatother
\makeatletter
\makeatother
\makeatletter
\@ifpackageloaded{caption}{}{\usepackage{caption}}
\@ifpackageloaded{subcaption}{}{\usepackage{subcaption}}
\makeatother
\usepackage{bookmark}
\IfFileExists{xurl.sty}{\usepackage{xurl}}{} % add URL line breaks if available
\urlstyle{same}
% fallback for those not using the hyperref driver hyperxmp:
\makeatletter
\@ifundefined{xmpquote}{\newcommand{\xmpquote}[1]{#1}}{}
\makeatother
\hypersetup{
  pdftitle={Práctica No.~6 -- Nutrición Mineral},
  pdflang={es},
  colorlinks=true,
  linkcolor={blue},
  filecolor={Maroon},
  citecolor={Blue},
  urlcolor={Blue},
  pdfcreator={LaTeX via pandoc}}


\title{\vspace{-1.5em}{\Huge \textbf{Práctica No.~6 -- Nutrición
Mineral}}\vspace{-1.5em}}
\date{}
\begin{document}
\maketitle


\setstretch{1.2}
\textbf{Fecha}: \hrulefill \hfill \textbf{Grupo}: \hrulefill

\textbf{Nombre:} \hrulefill

\textbf{Objetivos:}

\begin{itemize}
\tightlist
\item
  Comparar el crecimiento y desarrollo de plántulas cultivadas en
  condiciones de nutrición completa frente a aquellas sometidas a
  carencias de nutrientes.
\item
  Observar e identificar los síntomas fisiológicos de deficiencia que
  presentan las plantas cultivadas sin el aporte de elementos minerales
  esenciales.
\end{itemize}

\subsection{1. Materiales}\label{materiales}

\begin{itemize}
\tightlist
\item
  Plántulas o semillas germinadas de frijol (\emph{Phaseolus vulgaris})
  o maíz (\emph{Zea mays})
\item
  6 vasos plásticos o macetas pequeñas
\item
  Sustrato inerte (vermiculita)
\item
  Agua para riego (agua corriente almacenada en recipiente abierto)
\item
  Solución nutritiva (fertilizante Triple 15)
\item
  Papel absorbente o papel filtro
\item
  Regla graduada o papel milimetrado
\item
  Kit de disección
\end{itemize}

\subsection{2. Aspectos Teóricos e
Introducción}\label{aspectos-teuxf3ricos-e-introducciuxf3n}

La nutrición mineral es fundamental para el crecimiento, desarrollo y
rendimiento fisiológico de las plantas. Los nutrientes minerales se
dividen en \textbf{macronutrientes} (como nitrógeno, fósforo y potasio),
requeridos en grandes cantidades, y \textbf{micronutrientes} (como
hierro, zinc y cobre), necesarios en menores concentraciones pero
igualmente vitales. La deficiencia o exceso de estos elementos altera
procesos metabólicos esenciales (fotosíntesis, síntesis de proteínas y
clorofila) y se manifiesta visualmente mediante clorosis, necrosis o
enanismo.

Defina brevemente los siguientes conceptos:

\subsubsection{1. Macronutrientes vs.~Micronutrientes (funciones
principales):}\label{macronutrientes-vs.-micronutrientes-funciones-principales}

\hrulefill

\hrulefill

\hrulefill

\hrulefill

\hrulefill

\subsubsection{2. Síntomas de deficiencia nutricional (Clorosis,
Necrosis,
Enanismo):}\label{suxedntomas-de-deficiencia-nutricional-clorosis-necrosis-enanismo}

\hrulefill

\hrulefill

\hrulefill

\hrulefill

\hrulefill

\subsubsection{3. Sustrato inerte y su importancia en ensayos
nutricionales:}\label{sustrato-inerte-y-su-importancia-en-ensayos-nutricionales}

\hrulefill

\hrulefill

\hrulefill

\hrulefill

\hrulefill

\subsection{3. Procedimiento
Experimental}\label{procedimiento-experimental}

\subsubsection{A. Montaje e Inicio del Ensayo (Día 1 a Día
7)}\label{a.-montaje-e-inicio-del-ensayo-duxeda-1-a-duxeda-7}

\begin{enumerate}
\def\labelenumi{\arabic{enumi}.}
\tightlist
\item
  \textbf{Día 1:} Deje las semillas de fríjol o maíz en agua durante 24
  horas.
\item
  Siembre tres semillas por maceta en el sustrato inerte (vermiculita) y
  riegue con el agua para riego.
\item
  \textbf{Día 7:} Aplique los primeros tratamientos:

  \begin{itemize}
  \tightlist
  \item
    \textbf{Macetas 1, 2 y 3 (Tratamiento Nutritivo):} 10 mL de solución
    nutritiva cada una.
  \item
    \textbf{Macetas 4, 5 y 6 (Control Carencial):} 10 mL de agua de
    riego cada una.
  \end{itemize}
\end{enumerate}

\begin{center}
\begin{tikzpicture}
  % Esquema de Macetas
  \foreach \x/\txt in {0/M1 (Nut), 1.2/M2 (Nut), 2.4/M3 (Nut), 4.0/M4 (Ctrl), 5.2/M5 (Ctrl), 6.4/M6 (Ctrl)} {
    \draw[thick, fill=gray!10] (\x,0) -- (\x+0.2,-0.8) -- (\x+0.8,-0.8) -- (\x+1,0) -- cycle;
    \node[scale=0.6] at (\x+0.5,-0.4) {\txt};
  }
\end{tikzpicture}
\end{center}

\subsubsection{B. Seguimiento y Mantenimiento (Día 8 a Día
30)}\label{b.-seguimiento-y-mantenimiento-duxeda-8-a-duxeda-30}

\begin{enumerate}
\def\labelenumi{\arabic{enumi}.}
\tightlist
\item
  Mantenga las macetas bajo condiciones ambientales similares.
\item
  Riegue con agua de riego a medida que el sustrato se seque (se
  recomienda regar en horas de la tarde).
\item
  Cada 7 días registre la longitud del tallo, número de hojas y síntomas
  cualitativos de carencia.
\item
  \textbf{Día 23:} Aplique nuevamente 10 mL de solución nutritiva a las
  macetas 1, 2 y 3 y 10 mL de agua de riego a las macetas 4, 5 y 6.
\item
  \textbf{Día 30:} Lleve las plantas al laboratorio y extraigalas
  cuidadosamente sin dañar el sistema radicular para las mediciones
  finales.
\item
  Para cada tratamiento (\(n=9\) plántulas), calcule el
  \textbf{Promedio} (\(\bar{x}\)) y la \textbf{Desviación Estándar}
  (\(s\)):
\end{enumerate}

\[\bar{x} = \frac{\sum_{i=1}^{n} x_i}{n} \qquad s = \sqrt{\frac{\sum_{i=1}^{n} (x_i - \bar{x})^2}{n - 1}}\]

\subsubsection{Tabla 1. Seguimiento Temporal Parcial (Longitud Vástago
cm / N°
Hojas)}\label{tabla-1.-seguimiento-temporal-parcial-longitud-vuxe1stago-cm-n-hojas}

\begin{center}
\resizebox{\linewidth}{!}{%
\begin{tabular}{c|c|c|c|c|c}
\hline
\textbf{Maceta} & \textbf{Sem.} & \textbf{Día 7 (L | H)} & \textbf{Día 14 (L | H)} & \textbf{Día 21 (L | H)} & \textbf{Día 30 (L | H)} \\
\hline
 & A & \quad | \quad & \quad | \quad & \quad | \quad & \quad | \quad \\
M1 (Nut) & B & \quad | \quad & \quad | \quad & \quad | \quad & \quad | \quad \\
 & C & \quad | \quad & \quad | \quad & \quad | \quad & \quad | \quad \\
\hline
 & A & \quad | \quad & \quad | \quad & \quad | \quad & \quad | \quad \\
M2 (Nut) & B & \quad | \quad & \quad | \quad & \quad | \quad & \quad | \quad \\
 & C & \quad | \quad & \quad | \quad & \quad | \quad & \quad | \quad \\
\hline
 & A & \quad | \quad & \quad | \quad & \quad | \quad & \quad | \quad \\
M3 (Nut) & B & \quad | \quad & \quad | \quad & \quad | \quad & \quad | \quad \\
 & C & \quad | \quad & \quad | \quad & \quad | \quad & \quad | \quad \\
\hline
\multicolumn{2}{c|}{\textbf{Promedio ($\bar{x}$) Nut}} & & & & \\
\multicolumn{2}{c|}{\textbf{Desv. Est. ($s$) Nut}} & & & & \\
\hline\hline
 & A & \quad | \quad & \quad | \quad & \quad | \quad & \quad | \quad \\
M4 (Ctrl) & B & \quad | \quad & \quad | \quad & \quad | \quad & \quad | \quad \\
 & C & \quad | \quad & \quad | \quad & \quad | \quad & \quad | \quad \\
\hline
 & A & \quad | \quad & \quad | \quad & \quad | \quad & \quad | \quad \\
M5 (Ctrl) & B & \quad | \quad & \quad | \quad & \quad | \quad & \quad | \quad \\
 & C & \quad | \quad & \quad | \quad & \quad | \quad & \quad | \quad \\
\hline
 & A & \quad | \quad & \quad | \quad & \quad | \quad & \quad | \quad \\
M6 (Ctrl) & B & \quad | \quad & \quad | \quad & \quad | \quad & \quad | \quad \\
 & C & \quad | \quad & \quad | \quad & \quad | \quad & \quad | \quad \\
\hline
\multicolumn{2}{c|}{\textbf{Promedio ($\bar{x}$) Ctrl}} & & & & \\
\multicolumn{2}{c|}{\textbf{Desv. Est. ($s$) Ctrl}} & & & & \\
\hline
\end{tabular}%
}
\end{center}

Interprete los resultados parciales obtenidos durante la fase
experimental:

\hrulefill

\hrulefill

\hrulefill

\hrulefill

\hrulefill

\subsection{4. Registro de Resultados (Día
30)}\label{registro-de-resultados-duxeda-30}

Determine para cada planta la longitud de raíz y vástago (cm), el peso
fresco de raíz y vástago (g) y el número de hojas. Anote los cambios
cualitativos observados (coloración, morfología).

\subsubsection{Tabla 2. Mediciones Finales Destructivas (Día
30)}\label{tabla-2.-mediciones-finales-destructivas-duxeda-30}

\begin{center}
\resizebox{\linewidth}{!}{%
\begin{tabular}{c|c|c|c|c|p{7em}}
\hline
\textbf{Maceta} & \textbf{Sem.} & \textbf{L. Raíz (cm)} & \textbf{P.F. Raíz (g)} & \textbf{P.F. Vást. (g)} & \textbf{Síntomas / Observaciones} \\
\hline
 & A & & & & \\
M1 (Nut) & B & & & & \\
 & C & & & & \\
\hline
 & A & & & & \\
M2 (Nut) & B & & & & \\
 & C & & & & \\
\hline
 & A & & & & \\
M3 (Nut) & B & & & & \\
 & C & & & & \\
\hline
\multicolumn{2}{c|}{\textbf{Promedio ($\bar{x}$) Nut}} & & & & \multicolumn{1}{c}{---} \\
\multicolumn{2}{c|}{\textbf{Desv. Est. ($s$) Nut}} & & & & \multicolumn{1}{c}{---} \\
\hline\hline
 & A & & & & \\
M4 (Ctrl) & B & & & & \\
 & C & & & & \\
\hline
 & A & & & & \\
M5 (Ctrl) & B & & & & \\
 & C & & & & \\
\hline
 & A & & & & \\
M6 (Ctrl) & B & & & & \\
 & C & & & & \\
\hline
\multicolumn{2}{c|}{\textbf{Promedio ($\bar{x}$) Ctrl}} & & & & \multicolumn{1}{c}{---} \\
\multicolumn{2}{c|}{\textbf{Desv. Est. ($s$) Ctrl}} & & & & \multicolumn{1}{c}{---} \\
\hline
\end{tabular}%
}
\end{center}

Interprete los resultados finales:

\hrulefill

\hrulefill

\hrulefill

\hrulefill

\hrulefill

\subsection{5. Representación
Gráfica}\label{representaciuxf3n-gruxe1fica}

\subsubsection{Gráfica 1: Cinética de Crecimiento
Temporal}\label{gruxe1fica-1-cinuxe9tica-de-crecimiento-temporal}

Represente la cinética temporal de la longitud del vástago usando puntos
(promedio por tratamiento) con barras de error (desviación estándar) y
líneas.

\begin{center}
\begin{tikzpicture}[scale=0.85]
  % Grid ggplot2 theme_bw style
  \draw[gray!20, step=0.7cm] (0,0) grid (7,5);
  \draw[gray!50, step=1.4cm] (0,0) grid (7,5);
  \draw[thick] (0,0) rectangle (7,5);
  
  % Eje X (Tiempo)
  \node[below, font=\tiny] at (1.4,0) {Día 7};
  \node[below, font=\tiny] at (2.8,0) {Día 14};
  \node[below, font=\tiny] at (4.2,0) {Día 21};
  \node[below, font=\tiny] at (5.6,0) {Día 30};
  \node[below=8pt, font=\small] at (3.25,0) {\textbf{Tiempo (Días)}};
  
  % Eje Y (Longitud)
  \node[rotate=90, above=10pt, font=\small] at (0,2.25) {\textbf{Longitud Vástago (cm)}};
  
  % Convención de Leyenda
  \draw[fill=white, draw=gray] (0.2,3.3) rectangle (2.4,4.3);
  \draw[thick, blue] (0.3,4.0) -- (0.7,4.0); \node[right, font=\tiny] at (0.7,4.0) {Nutritivo};
  \draw[thick, red, dashed] (0.3,3.6) -- (0.7,3.6); \node[right, font=\tiny] at (0.7,3.6) {Control};
\end{tikzpicture}
\end{center}

Interprete los resultados observados:

\hrulefill

\hrulefill

\hrulefill

\hrulefill

\hrulefill

\subsubsection{Gráfica 2: Biomasa Fresca
Final}\label{gruxe1fica-2-biomasa-fresca-final}

Trace un gráfico de cajas y bigotes (\emph{boxplot}) para el Peso Fresco
del Vástago por tratamiento.

\begin{center}
\begin{tikzpicture}[scale=0.85]
  % Grid ggplot2 theme_bw style
  \draw[gray!20, step=0.7cm] (0,0) grid (7,5);
  \draw[gray!50, step=1.4cm] (0,0) grid (7,5);
  \draw[thick] (0,0) rectangle (7,5);
  
  % Eje X (Tratamientos)
  \node[below, font=\small] at (2.1,0) {\textbf{Nutritivo}};
  \node[below, font=\small] at (4.9,0) {\textbf{Control}};
  \node[below=12pt, font=\small] at (3.25,0) {\textbf{Tratamiento Nutricional}};
  
  % Eje Y (Peso Fresco)
  \node[rotate=90, above=10pt, font=\small] at (0,2.25) {\textbf{Peso Fresco Vástago (g)}};
\end{tikzpicture}
\end{center}

Interprete los resultados observados:

\hrulefill

\hrulefill

\hrulefill

\hrulefill

\hrulefill

\subsection{6. Actividades
Complementarias}\label{actividades-complementarias}

\begin{enumerate}
\def\labelenumi{\arabic{enumi}.}
\tightlist
\item
  ¿De qué manera influyen las condiciones del suelo (pH, textura y
  composición orgánica) en la disponibilidad y absorción de nutrientes
  minerales, y qué estrategias pueden emplear las plantas para optimizar
  su nutrición en suelos desfavorables?
\end{enumerate}

\hrulefill

\hrulefill

\hrulefill

\hrulefill

\begin{enumerate}
\def\labelenumi{\arabic{enumi}.}
\setcounter{enumi}{1}
\tightlist
\item
  Analice el papel de los micronutrientes en la regulación de procesos
  metabólicos esenciales (como la fotosíntesis y la síntesis de
  clorofila) y explique cómo se pueden diagnosticar y corregir sus
  deficiencias en un cultivo agrícola.
\end{enumerate}

\hrulefill

\hrulefill

\hrulefill

\hrulefill

\begin{enumerate}
\def\labelenumi{\arabic{enumi}.}
\setcounter{enumi}{2}
\tightlist
\item
  ¿Cómo interactúa la nutrición mineral con la actividad de la
  microbiota del suelo y de qué manera esta sinergia contribuye a la
  salud y productividad de las plantas?
\end{enumerate}

\hrulefill

\hrulefill

\hrulefill

\hrulefill

\subsection{7. Conclusión}\label{conclusiuxf3n}

Compare el crecimiento y desarrollo de las plántulas entre tratamientos
e identifique las implicaciones agronómicas de los resultados obtenidos:

\hrulefill

\hrulefill

\hrulefill

\hrulefill

\subsection{8. Bibliografía}\label{bibliografuxeda}

\begin{itemize}
\item
  Azcón-Bieto, J., \& Talón, M. (2008). \emph{Fundamentos de Fisiología
  Vegetal} (2.ª ed.). McGraw-Hill Interamericana.
\item
  Melgarejo, L. M. (Ed.). (2010). \emph{Experimentos en fisiología
  vegetal}. Universidad Nacional de Colombia.
\item
  Taiz, L., \& Zeiger, E. (2002). \emph{Plant Physiology} (3rd ed.).
  Sinauer Associates Ltd.
\end{itemize}




\end{document}
